% Generated draft: requires mathematical and source review.
\documentclass[11pt]{article}
\usepackage[margin=25mm]{geometry}
\usepackage{fontspec}
\setmainfont{Latin Modern Roman}
\usepackage{amsmath,amssymb,mathtools,unicode-math}
\setmathfont{Latin Modern Math}
\usepackage{enumitem,needspace,xurl,hyperref}
\setlength{\parindent}{0pt}
\setlength{\parskip}{5pt}
\setlength{\emergencystretch}{3em}
\newcommand{\N}{\mathbb{N}}
\newcommand{\Z}{\mathbb{Z}}
\newcommand{\Q}{\mathbb{Q}}
\newcommand{\R}{\mathbb{R}}
\newcommand{\C}{\mathbb{C}}
\newcommand{\F}{\mathbb{F}}
\newcommand{\A}{\mathbb{A}}
\newcommand{\PP}{\mathbb{P}}
\newcommand{\E}{\mathbb{E}}
\newcommand{\eps}{\varepsilon}
\newcommand{\dd}{\,\mathrm{d}}
\begin{document}
\section*{UnsolvedMath ID: 6700069}
\textbf{Source problem code:} catalogue entry
\subsection{Definitions and mathematical statement}
Let \(S\) be a Riemannian manifold homeomorphic to
\[
 \#_{i=1}^{20}(S^2\times S^2).
\]
In particular it is closed and orientable; choose an orientation. Let \(B^4(R)\subset\R^4\) be the Euclidean ball of radius \(R\), and give \(S\) its specified Riemannian distance. A map \(f_R:B^4(R)\to S\) is 1-Lipschitz if \(d_S(f_R(x),f_R(y))\leq|x-y|\) for all \(x,y\).

Choose a cocycle \(h\) representing the top dimensional cohomology orientation class \([S]\in H^4(S;\R)\). In the source's differential-form version, take a smooth closed 4-form \(h\) normalized by \(\int_S h=1\), and interpret
\[
 h(f_R)=\int_{B^4(R)}f_R^*h,
\]
using the almost everywhere differential of a Lipschitz map. The question is whether one can find such maps for radii tending to infinity, with a positive constant \(c=c(S)\) independent of \(R\), satisfying
\[
 \int_{B^4(R)}f_R^*h\ \geq\ cR^4.
\]
The integral is signed; no absolute value is inserted. The form is fixed independently of \(R\), and there is no stated boundary condition on \(f_R\).

The source does not define \(h(f_R)\) for an arbitrary cocycle. Since the domain ball has boundary, this is not automatically a pairing of an absolute fundamental cycle with the cohomology class; a cocycle evaluation convention or the differential-form interpretation is needed. The notation \(R\to\infty\) also leaves open whether maps are required for every sufficiently large radius or merely along an unbounded sequence. Both readings retain a single uniform positive constant.
\subsection{Short English statement}
Can large Euclidean four-balls map 1-Lipschitzly to the connected sum of twenty copies of \(S^2\times S^2\), with a fixed top dimensional cocycle evaluating at least proportionally to \(R^4\)? The differential-form version uses a signed pullback integral and imposes no boundary condition.
\subsection{Sources}
\begin{itemize}
\item Local UnsolvedMath registry, permanent ID 6700069, source code catalogue entry; Misha Gromov, 2017; update 2026-08-06.
\item Gromov, \emph{101 Questions, Problems and Conjectures around Scalar Curvature} (2017), printed tag [?73], occurrence 72, PDF page 72, subquestion [c]: \url{https://www.ihes.fr/~gromov/wp-content/uploads/2018/08/101-problemsOct1-2017.pdf}. Registry access date 2026-07-29; normalized PDF extraction; public author-hosted incomplete PDF; present status and reuse terms require review.
\item The supplied summary describes a Lipschitz map and cohomological growth question but supplies no precise further reference or resolution.
\end{itemize}
\end{document}
